% Stand 18. Juli 2026
\subsubsection{Analyse der direktionalen Variogramme}
\label{KapitelDatenauswertung_Ergebnisse_Variogramme}
Analog zu der in Kapitel \ref{KapitelTheorieVariogramm} 
eingeführten ungerichteten empirischen Semivarianz 
erfordert auch die Untersuchung der richtungsabhängigen 
empirischen  Semivarianz eine Vorauswahl hinsichtlich 
Schätzfunktion und Kovarianz-Modell.\\
Ein guter Startpunkt für die zu treffende Wahl 
sind die Schätzfunktion nach Cressie-Hawkins
und das Matérn Kovarianz-Modell: ersteres, da es sich in
Voruntersuchungen als  besonders robust erwiesen 
hat; letzteres, da dessen funktionale Beschreibung durch den
zusätzlichen Parameter $\nu$ ein hohe Flexibilität besitzt.\\
Die Anpassung (\textit{engl.:} \glqq fit\grqq ) der Schätzungen an das Parameter-Tupel 
(Sill, Range, Nugget) der theoretischen 
Kovarianz-Funktion erfolgt mithilfe eines 
Optimierungs-Verfahren\footnote{Es wird der 
\glqq \textbf{T}rust \textbf{R}egion
\textbf{R}eflective\grqq\ Optimierungs-Algorithmus verwendet,
welches oft für die nicht-lineare Kurvenanpassung zum Einsatz kommt.}, dem ein beschränkendes Intervall für 
den anzupassenden Parameter übergeben wird.\\
Für den Range Parameter wird dieses Intervall mit
$\mathcal{J}^{(R)}_{fit}=
[0\,m,\,R_{max}]$ festgelegt, wobei der \textit{obere} 
Wert $R_{max}$ durch die maximale paarweise
euklidische Distanz der
betrachteten Sensoren vorgegeben ist 
(z.\,B.\,wäre in 
Abb.\,(\ref{Abb_Datenauswertung_VerteilungDistanzenApr2026})
$R_{max}=22\,km$).\\

Aus rein mathematischer Perspektive ist die Beschränkung des
Range Parameters durch $R_{max}$ \textit{nicht} notwendig.
Die Begründung hierfür hat eher numerische Gründe und berührt Fragen 
des wissenschaftlichen Stils:
\begin{itemize}[topsep=2pt,parsep=1pt,itemsep=0.5pt]
\item[a) ] falls der \glqq wahre\grqq\ Range die beobachteten Entfernungen 
überschreitet, kann es dazu kommen, dass viele Modelle mit großen 
geschätzten Werten für den Range nahezu identische Residuen 
aufweisen: das Optimierungs-Problem ist in diesem Falle schlecht konditioniert.
\item[b) ] Wenn der Sill innerhalb der beobachteten Entfernungen nicht erreicht wird,
enthalten die Daten nicht genügend Information, um den \glqq wahren\grqq\ Range 
zu schätzen. Die Angabe der größten beobachteten Distanz $R_{max}$ als Schätzung für 
den Range erscheint unter diesem Blickwinkel \glqq aufrichtiger\grqq als
der Rückgabewert einer unbeschränkten und möglicherweise unplausiblen
 numerischen Anpassung.
\end{itemize}

\vspace*{5pt}

Mithilfe dieser Vorauswahl und mit einer Winkel-Schrittweite von 
$20^\circ$ werden die direktionalen Semivarianzen geschätzt 
und an das festgelegte Kovarianz-Modell angepasst.\\

\begin{figure}[H] 
\centering
\includegraphics[scale = 0.45]{Datenauswertung_P25MessungenApr2026DirektionaleSemivarianzBeispielFit340} 
\caption{\small{Anpassung der Kovarianz-Funktion 
(\textcolor{NavyBlue}{\rule[1mm]{3mm}{0.5mm}}) an die
direktionale ($\vartheta=340^\circ$) empirische Semivarianz 
(\textcolor{Sepia}{$\bullet$});
oberhalb der empirischen Semivarianz sind die Umfänge 
des Lags vermerkt}}
\label{Abb_Datenauswertung_P25MessungenApr2026DirektionaleSemivarianzBeispielFit340}
\end{figure} 

\begin{figure}[H] 
\centering
\includegraphics[scale = 0.45]{Datenauswertung_P25MessungenApr2026DirektionaleSemivarianzBeispielFit140} 
\caption{\small{Anpassung der Kovarianz-Funktion 
(\textcolor{NavyBlue}{\rule[1mm]{3mm}{0.5mm}}) an die
direktionale ($\vartheta=140^\circ$) empirische Semivarianz 
(\textcolor{Sepia}{$\bullet$});
oberhalb der empirischen Semivarianz sind die Umfänge 
des Lags vermerkt}}
\label{Abb_Datenauswertung_P25MessungenApr2026DirektionaleSemivarianzBeispielFit140}
\end{figure} 

Die 
Abb.\,(\ref{Abb_Datenauswertung_P25MessungenApr2026DirektionaleSemivarianzBeispielFit340}) 
und 
(\ref{Abb_Datenauswertung_P25MessungenApr2026DirektionaleSemivarianzBeispielFit140})
zeigen zwei typische Ergebnisse dieser Anpassung der
Matérn Kova\-ri\-anz-Funktion an die richtungsabhängig geschätzten Semivarianzen
basierend auf den gemittelten 
$\mathrm{P}_{2.5}$-Messwerten für April 2026.\\
Drei Dinge fallen an diesen Abbildungen auf:
\begin{itemize}[topsep=2pt,parsep=1pt,itemsep=0.5pt]
\item[1)] Die größten Distanzklassen (\textit{Lags}) 
der Semivarianz-Schätzung (\textcolor{Sepia}{$\bullet$}) enthalten
erheblich weniger Distanzpaare als die kleinsten 
(die Anzahl der Paare in den Lags ist in Klammern
 oberhalb der Symbole vermerkt). Dies ist deswegen 
erwähnenswert, weil gerade die größten Lags maßgeblich zur
 Bestimmung des Range des theoretischen Modells beitragen 
 und geringe Umfänge der Distanzklassen die Stabilität der Schätzung 
 nachteilig beeinflussen.
\item[ 2) ] Die empirische Semivarianz erreicht in keiner der 
beiden Winkelrichtungen das Sill-Plateau, wie es bei hinreichend 
großer Wahl des Lags stattfinden sollte. 
\item[3) ] Als Folge davon war es in
Abb.\,(\ref{Abb_Datenauswertung_P25MessungenApr2026DirektionaleSemivarianzBeispielFit340})
für den Algorithmus optimal, den Range auf den größten, möglichen Wert 
von $R_{max}$ des beschränkenden Intervalls $\mathcal{J}^{\mathrm{(R)}}_{fit}$ 
zu setzen. Bei der Anpassung in
Abb.\,(\ref{Abb_Datenauswertung_P25MessungenApr2026DirektionaleSemivarianzBeispielFit140})
wurde der Range etwa in die Mitte des Intervalls
$\mathcal{J}^{\mathrm{(R)}}_{fit}$ gesetzt, so dass links und rechts, 
bzw.\, ober- und unterhalb der angepassten Werte für Range und Sill etwa 
gleich viele der geschätzten empirischen Semivarianzen liegen.
\end{itemize}

Diese Tendenzen gelten auch für andere, mögliche 
Kombinationen  von Schätzfunktion der Semivarianz und 
Kovarianz-Modell, weshalb die getroffenen
 Beobachtungen mit großer Wahrscheinlichkeit nicht durch eine
ungünstige Wahl von Schätzer und theoretischem Modell verursacht
sind, sondern aus der versuchsbedingten, räumlichen Verteilung 
der Sensoren resultieren. 

\begin{figure}[H] 
\centering
\includegraphics[scale = 0.45]{Datenauswertung_P25_Apr2026_VariogrammKarte_Range_Matern} 
\caption{\small{Direktionale Variogramm-Karte für den 
Range basierend auf den gemittelten 
$\mathrm{P}_{2.5}$- Messwerten für April 2026}}
\label{Abb_Datenauswertung_P25_Apr2026_VariogrammKarte_Range_Matern}
\end{figure} 

In Abb.\,(\ref{Abb_Datenauswertung_P25_Apr2026_VariogrammKarte_Range_Matern}) 
ist die direktionale Variogramm-Karte wiedergegeben, die 
sich aus allen 18 richtungsabhängigen Schätzungen der Semivarianz zusammensetzt.
Diese Art der Darstellung, bei der die 
Anpassungen an theoretische Kovarianz-Modelle für
richtungsabhängige Schätzungen der Semivarianz dargestellt
sind, orientiert sich an den Rosen-Diagrammen und wird 
\,z.\,B.\,in \cite[Kap.\,8.1]{Agterberg2014} verwendet.\\
Die Winkelangaben der 
Abb.\,(\ref{Abb_Datenauswertung_P25_Apr2026_VariogrammKarte_Range_Matern}) 
beziehen sich auf die Drehung im Uhrzeigersinn
gegen die Ost-West-Achse. Die $360^\circ / 20^\circ=18$ blauen
Punkte (\textcolor{CornflowerBlue}{$\bullet$}) kennzeichnen 
die Parame\-ter-Anpassungen der geschätzten Semivarianz in die 
angegeben Winkelrichtung. Die Parameter-Anpassungen sind durch
eine Linie miteinander verbunden und der sich ergebende
Polygonzug ist bläulich 
(\textcolor{CornflowerBlue}{\rule[-0.3mm]{3mm}{3mm}}) 
eingefärbt. Hierdurch lassen sich eventuelle
Abweichungen von der Rotations-Symmetrie 
(was einem Kreis entspräche) besonders leicht 
erkennen. \\
Bei Betrachtung der 
Abb.\,(\ref{Abb_Datenauswertung_P25_Apr2026_VariogrammKarte_Range_Matern})
wird eine Besonderheit sichtbar, die bereits im Kontext 
der Diskussion der 
Abb.\,(\ref{Abb_Datenauswertung_P25MessungenApr2026DirektionaleSemivarianzBeispielFit340}) 
und 
(\ref{Abb_Datenauswertung_P25MessungenApr2026DirektionaleSemivarianzBeispielFit140})
aufgezeigt wurde: bei der Anpassung des Parameters Range 
\glqq klebt\grqq\ der Algorithmus oft an der oberen Schranke fest, 
sodass der in diesen Fällen durch die Anpassung ermittelte Wert für
 den Range mit hoher Wahrscheinlichkeit ein Artefakt 
der Optimierung aber nicht den \glqq wahren\grqq\ Wert widerspiegelt.\\

Dieses Verhalten wurde für alle drei der verfügbaren, 
monatlich gemittelten Daten beobachtet, sowohl für die Messungen der
Konzentration von $\mathrm{P}_{2.5}$ wie auch von 
$\mathrm{P}_{10}$. Setzt man den oberen Wert des Intervalls 
$\mathcal{J}^{\mathrm{(R)}}_{fit}$ für die Range Anpassung 
auf den unrealistischen Wert von $100\,km$ zeigt sich dasselbe
Verhalten bei Betrachtung der Variogramm-Karte: der 
Optimierungs-Algorithmus setzt den Range-Parameter 
auf diesen maximal zulässigen Wert fest. 
Ein Beispiel der Anpassung der Kovarianz-Funktion an die
direktionale empirische Semivarianz sowie die direktionale
Variogramm-Karte unter dieser Nebenbedingung  
befinden sich im Anhang 
(Abb.\,(\ref{Abb_Datenauswertung_P25MessungenApr2026DirektionaleSemivarianzBeispielFit340_NoConstraints})
und
Abb.\,(\ref{Abb_Datenauswertung_P25MessungenApr2026VariogrammRoseRange_NoConstraints})
). 

Abb.\,(\ref{Abb_Datenauswertung_P25_Apr2026_VariogrammKarte_Range_Matern})
allein lässt somit keine Aussage über eine mögliche Anisotropie
oder der darauf aufbauenden Unterscheidung geometrische versus zonale Anisotropie zu.\\

\begin{figure}[H] 
\centering
\includegraphics[scale = 0.45]{Datenauswertung_P25_Nov2025_VariogrammKarte_Sill_Matern} 
\caption{\small{Direktionale Variogramm-Karte für den 
Sill basierend auf den gemittelten 
$\mathrm{P}_{2.5}$- Messwerten für November 2025}}
\label{Abb_Datenauswertung_P25_Nov2025_VariogrammKarte_Sill_Matern}
\end{figure} 

Interessanterweise zeigt die direktionale Variogramm-Karte 
für den Sill\footnote{Zur Erinnerung: Der Sill bezeichnet das Niveau, auf welches die 
Semivarianz für große Entfernungen 
konvergiert. Insbesondere beinhaltet der Sill den ggf.\,vorhandenen Nugget-Effekt.} eine ausgeprägtere Bevorzugung einer Richtung auf, wie aus der
Abb.\,(\ref{Abb_Datenauswertung_P25_Nov2025_VariogrammKarte_Sill_Matern}) 
hervorgeht.

In dieser Abbildung wurde zusätzlich die Achse eingezeichnet, 
entlang der der größte Wert für den Sill
ermittelt wurde (\textcolor{RubineRed}{\rule[0.4mm]{3mm}{1mm}}), sowie
senkrecht dazu eine Achse derselben Länge. Durch diese 
Konstruktion soll die visuelle Abschätzung 
der Asymmetrie erleichtert werden. 
So ist der Wert des Sill entlang der Achse 
$40^\circ \leftrightarrow 220^\circ$ in
Abb.\,(\ref{Abb_Datenauswertung_P25_Nov2025_VariogrammKarte_Sill_Matern})
um den Faktor 2-3 höher als in der senkrecht dazu
stehenden Richtung.\\
Diese Asymmetrie fällt für die Daten der Monate April 2026 und 
Juli 2025 deutlich schwächer aus; die zugehörige Abbildung 
befindet sich im Anhang 
(vgl.\,Abb.\,(\ref{Abb_Datenauswertung_P25_Apr2026Jul2025_VariogrammKarte_Sill_Matern})). \\

Nun stellen direktionale Variogramm-Karten des Sill 
keine Standardmethode zur 
Detektion einer möglichen Anisotropie dar,
welche sich ja vorrangig in der Richtungsabhängigkeit des
Range offenbart. Andererseits charakterisieren Sill und
Range denselben 
ausgezeichneten Punkt in der Variogramm-Analyse: 
sie bezeichnen die $y$- und $x$-Koordinate des Punktes,
an dem das Variogramm in das Plateau der Unkorreliertheit 
einmündet.\\ 
Diese Korrelation lässt sich 
anschaulich verstehen: 
Für viele der theoretischen Kovarianz-Modelle hat die
isotrope Semivarianz die Form (vgl.\,Gl.\,(\ref{Gl_Semivarianz_isotroperFall}))
\begin{equation}
\gamma(r) = c_0 + \sigma^2 f\left( \frac{r}{a}\right),
\end{equation} mit
\begin{itemize}[topsep=2pt,parsep=1pt,itemsep=0.5pt]
\item[$\circ$] dem Nugget $c_0$,
\item[$\circ$] dem Sill $\sigma^2$,
\item[$\circ$] dem Range $a$ und
\item[$\circ$] einer Funktion $f$ mit $f(0) = 0$ und 
$\displaystyle \lim_{r \longrightarrow \infty} f(r)=1$.
\end{itemize}
Damit folgt
\begin{equation}
\gamma'(r) := \frac{\mathrm{d}\gamma(r)}{\mathrm{d}r}=
\frac{\sigma^2}{a}\cdot f'(r)
\end{equation}
Wenn der Sill noch nicht klar erreicht wird und somit die Daten im Wesentlichen nur den 
anfänglichen Anstiegsbereich gut abbilden, dann bestimmt die Anpassung vorrangig 
das Verhältnis $\frac{\sigma^2}{a}$ aber nicht die Einzelgrößen $\sigma^2$ und $a$. 
Eine Erhöhung von $a$ muss somit zu einem Anwachsen von $\sigma^2$ führen, damit die 
Kurvenanpassung weiterhin zu den beobachteten Punkten passt. 
Das führt zu einer positiven Korrelation der beiden Größen Sill und Range. 
Theoretisch fundiert lässt sich weitergehend sogar zeigen, dass unter bestimmten 
Voraussetzungen Sill und Range \textit{nicht} konsistent unabhängig voneinander geschätzt werden 
können (vgl.\,\cite{Zhang2004}). 


Somit \textit{könnte} ein winkelabhängige Asymmetrie des Sill 
auf eine Winkelabhängigkeit des Range und somit indirekt auf eine 
Anisotropie hindeuten. \\

Da die Analyse der Variogramm-Karten \textit{indirekt} Hinweise auf mögliche 
Anisotropie liefert, wird im folgenden 
Kapitel \ref{KapitelDatenauswertung_Ergebnisse_Kriging} untersucht, 
ob sich diese Indizien auch in der Vorhersagequalität des Kriging 
widerspiegeln. Hierzu wurden anisotrope Kovarianz-Modelle konstruiert 
und mittels einer Variante des Cross Validation bewertet.




 
